The domain $\Omega=\{(x,y)\in \mathbb{R}^2 \mid 0<x<1\wedge 0<y<1\}$
in the coordinate system of figure~\ref{40000033:domain} is a rhombus.
\begin{figure}[h]
\centering
\begin{tikzpicture}[>=latex,thick,scale=4]
\fill[color=red!20] (0,0) -- (1,0) -- +(60:1) -- (60:1) -- cycle;
\fill[color=gray!20,opacity=0.5] (0,0) -- (0:0.3) arc(0:60:0.3) -- cycle;
\draw[color=red,line width=1.4pt] (0,0) -- (1,0) -- +(60:1) -- (60:1) -- cycle;
\node[color=red] at ($0.5*(1,0)+0.5*(60:1)$) {$\Omega$};
\draw[->] (-0.1,0) -- (1.6,0) coordinate[label={$x$}];
\draw[->] (-120:0.1) -- (60:1.6) coordinate[label={$y$}];
\node at (30:0.2) {$\frac{\pi}3$};
\draw (1,-0.02) -- (1,0.02);
\node at (1,0) [below] {$1$};
\draw ($1*(60:1)+1*(-30:0.02)$) -- +(150:0.04);
\node at (60:1) [above left] {$1$};
\end{tikzpicture}
\caption{Domain for problem~\ref{40000033}
\label{40000033:domain}}
\end{figure}
Determine the values $\lambda$, for which the
partial differential equation 
\[
\biggl(
\frac{\partial^2}{\partial x^2}
+
\frac{\partial^2}{\partial y^2}
\biggr) u = \lambda u
\]
with homogenous boundary conditions has a solution.

\begin{loesung}
We write the operator as
\[
L
=
\frac{\partial^2}{\partial x^2}
+
\frac{\partial^2}{\partial y^2}.
\]
We write the solution $u(x,y)$ as a product $u(x,y)=X(x)Y(y)$.
Inserting this into the differential equations gives
\[
L u
=
X''(x)Y(y) + X(x) Y''(y)
=
\lambda X(x) Y(y).
\]
Dividing by $X(x) Y(y)$ gives
\[
\frac{X''(x)}{X(x)}
+
\frac{Y''(y)}{Y(y)}
=
\lambda
\quad
\Rightarrow
\quad
\frac{X''(x)}{X(x)}
=
\lambda
-
\frac{Y''(y)}{Y(y)}.
\]
As the left side only depends on $x$ and the right only on $y$, both
sides need to be constant, we call this constant $-\mu^2$, giving
the ordinary differential equations
\begin{align*}
X''(x)
&=
-\mu^2 X(x)
&
Y''(x)
&=
(\lambda + \mu^2)Y(x)
\intertext{with solutions}
X(x)
&=
\sin\mu x
&
Y(y)
&=
\sin \!\sqrt{\lambda+\mu^2} y
\end{align*}
From the homogeneous boundary conditions we conclude that
$0=X(1) = \sin\mu$ which implies $\mu = \pi k$ with $k\in\mathbb{Z}$.
Similarly,
$\pi l=\!\sqrt{\lambda+\mu^2}=\!\sqrt{\lambda+\pi^2 k^2}\in\mathbb{N}$.
This means that
\[
\lambda = -\pi^2 (k^2 +l^2).
\]
Put together, the solutions are
\[
u_{k,l}(x,y)
=
\sin(\pi k x)
\sin(\pi l y)
\]
with $k,l\in \mathbb{N}$.
\end{loesung}

\begin{diskussion}
Note that the operator $L$ is not the Laplace operator,
The correct form of the Laplace operator can be derived as follows.
Rectangular coordinates $(\xi,\eta)$ can be expressed by $x$ and $y$
using
\[
\renewcommand{\arraycolsep}{3pt}
\begin{array}{rcrcr}
\xi  &=& x &+& \frac12 y          \\
\eta &=&   & & \frac{\!\sqrt{3}}2 y
\end{array}
\qquad\Leftrightarrow\qquad
\begin{array}{rcrcr}
x &=& \xi &-& \frac{1}{\!\sqrt{3}} \eta\\
y &=&     & & \frac{2}{\!\sqrt{3}} \eta
\end{array}
\]
This equations imply the partial derivatives
\begin{align*}
\frac{\partial \xi}{\partial x}  &= 1 
&
\frac{\partial x}{\partial \xi}  &= 1
\\
\frac{\partial \xi}{\partial y}  &= \frac12
&
\frac{\partial x}{\partial \eta} &= -\frac{1}{\!\sqrt{3}}
\\
\frac{\partial \eta}{\partial x} &= 0
&
\frac{\partial y}{\partial \xi}  &= 0
\\
\frac{\partial \eta}{\partial y} &= \frac{\!\sqrt{3}}2
&
\frac{\partial y}{\partial \eta} &= \frac{2}{\!\sqrt{3}}
\end{align*}
This can be used to express the first and second derivatives with
respect to the rectangular coordinates in terms of derivatives with
respect to $x$ and $y$.
They are
\begin{align*}
\frac{\partial}{\partial \xi}
&=
\frac{\partial x}{\partial \xi}\frac{\partial}{\partial x}
+
\frac{\partial y}{\partial\xi}\frac{\partial}{\partial y}
=
\frac{\partial}{\partial x}
&
\frac{\partial}{\partial\eta}
&=
\frac{\partial x}{\partial\eta}\frac{\partial}{\partial x}
+
\frac{\partial y}{\partial\eta}\frac{\partial}{\partial y}
=
-\frac{1}{\!\sqrt{3}}\frac{\partial}{\partial x}
+\frac{2}{\!\sqrt{3}}\frac{\partial}{\partial y}
\end{align*}
and
\[
\Delta
=
\frac{\partial^2}{\partial\xi^2}
+
\frac{\partial^2}{\partial\eta^2}
=
\frac{\partial^2}{\partial x^2}
+\frac13\frac{\partial^2}{\partial x^2}
-\frac{4}{3}\frac{\partial^2}{\partial x\,\partial y}
+\frac{4}{3}\frac{\partial^2}{\partial y^2}
=
\frac43\biggl(
\frac{\partial^2}{\partial x^2}
-\frac{\partial^2}{\partial x\,\partial y} + \frac{\partial^2}{\partial y^2}
\biggr)
\]
This shows that $L$ ist not $\Delta$.
\end{diskussion}

\begin{bewertung}
Product ansatz ({\bf A}) 1 point,
separation of variables ({\bf S}) 1 point,
two separated differential equations ({\bf D}) 1 point,
solutions of the equations ({\bf S}) 1 point,
restriction of the value of $\mu$ from boundary conditions ({\bf B}) 1 point,
possible values for $\lambda$ ({\bf L}) 1 point.
\end{bewertung}

